%% proposta_voo_bistatico.tex
%% Proposta técnica: otimização do plano de voo do receptor em SAR biestático
%% (transmissor em hélice cônica, receptor com posição otimizada),
%% para os cenários de solo homogêneo e solo estratificado (multicamada).
%% Idioma: Português
%% Compilar: pdflatex proposta_voo_bistatico.tex (duas passagens)

\documentclass[11pt,a4paper]{article}

% ─── pacotes ────────────────────────────────────────────────────────────────
\usepackage[utf8]{inputenc}
\usepackage[T1]{fontenc}
\usepackage[portuguese]{babel}
\usepackage{lmodern}
\usepackage{geometry}
\geometry{left=2.8cm, right=2.5cm, top=2.8cm, bottom=2.8cm}

\usepackage{amsmath,amssymb,amsfonts}
\usepackage{mathtools}
\usepackage{bm}
\usepackage{siunitx}
\sisetup{output-decimal-marker={,}}

\usepackage{graphicx}
\usepackage{float}
\usepackage{caption}
\usepackage[table]{xcolor}
\usepackage{booktabs}
\usepackage{array}
\usepackage{multirow}
\usepackage{longtable}
\usepackage{enumitem}
\usepackage{setspace}
\onehalfspacing

\usepackage{tikz}
\usetikzlibrary{arrows.meta,calc,patterns,positioning}

\usepackage{tcolorbox}
\tcbuselibrary{skins,breakable}

\usepackage[unicode,colorlinks=true,
  linkcolor=blue!60!black,
  citecolor=green!50!black,
  urlcolor=cyan!60!black,
  pdftitle={Proposta de Voo Bistatico para SAR de Subsuperficie},
  pdfauthor={David Atencia}]{hyperref}

\usepackage{fancyhdr}
\setlength{\headheight}{15pt}
\pagestyle{fancy}
\fancyhf{}
\fancyhead[L]{\small\nouppercase{\leftmark}}
\fancyhead[R]{\small\thepage}
\renewcommand{\headrulewidth}{0.4pt}

% ─── caixas de destaque ───────────────────────────────────────────────────────
\newtcolorbox{resultado}[1]{enhanced, breakable, colback=yellow!8, colframe=orange!70!black,
  fonttitle=\bfseries, title={#1}}
\newtcolorbox{definicao}[1]{enhanced, breakable, colback=blue!5, colframe=blue!50!black,
  fonttitle=\bfseries, title={#1}}
\newtcolorbox{nota}[1]{enhanced, breakable, colback=green!5, colframe=green!50!black,
  fonttitle=\bfseries\small, title={#1}}

% ─── macros ────────────────────────────────────────────────────────────────
\newcommand{\RT}{R_T}
\newcommand{\RR}{R_R}
\newcommand{\thetaB}{\theta_B}
\newcommand{\nuno}{n_1}
\newcommand{\ndos}{n_2}
\newcommand{\Prcv}{P_r}
\newcommand{\Pt}{P_t}
\newcommand{\Jcost}{J}
\newcommand{\alphaR}{\alpha}
\newcommand{\dxy}{\delta_{xy}}
\newcommand{\dz}{\delta_z}

% ─── capa moderna ─────────────────────────────────────────────────────────────
\definecolor{covernavy}{HTML}{102A43}
\definecolor{coverteal}{HTML}{1B998B}
\definecolor{coversub}{HTML}{627D98}
\newcommand{\coverfont}{\sffamily}

\begin{document}

\begin{titlepage}
\newgeometry{left=0pt,right=0pt,top=0pt,bottom=0pt}
\thispagestyle{empty}
\begin{tikzpicture}[remember picture, overlay]

  % ── faixa lateral navy ──────────────────────────────────────────────────
  \fill[covernavy] (current page.north west) rectangle ([xshift=5.4cm]current page.south west);
  \fill[coverteal] ([xshift=5.4cm]current page.north west) rectangle ([xshift=5.5cm]current page.south west);

  % ── anéis concêntricos decorativos (motivo radar), recortados à faixa ──
  \begin{scope}
    \clip (current page.north west) rectangle ([xshift=5.4cm]current page.south west);
    \foreach \r in {1.6,3.4,5.2,7.0,8.8,10.6}{
      \draw[coverteal!55, line width=0.7pt]
        ([yshift=-6.2cm,xshift=1.9cm]current page.north west) circle (\r cm);
    }
    \fill[coverteal!85] ([yshift=-6.2cm,xshift=1.9cm]current page.north west) circle (2.2pt);
  \end{scope}

  % ── rótulo vertical no rodapé da faixa ──────────────────────────────────
  \node[rotate=90, anchor=west]
    at ([xshift=2.7cm,yshift=2.0cm]current page.south west)
    {\coverfont\bfseries\fontsize{9}{9}\selectfont\color{coverteal!90!white}
     SAR BIEST\'ATICO \;\textbullet\; TOMOGRAFIA DE SUBSUPERF\'ICIE};

  % ── kicker ───────────────────────────────────────────────────────────────
  \node[anchor=north west, text width=12.6cm, align=left]
    at ([xshift=6.6cm,yshift=-9.6cm]current page.north west)
    {\coverfont\bfseries\color{coverteal} \large PROPOSTA T\'ECNICA};

  % ── título principal ────────────────────────────────────────────────────
  \node[anchor=north west, text width=12.6cm, align=left]
    at ([xshift=6.6cm,yshift=-11.1cm]current page.north west)
    {\coverfont\bfseries\fontsize{25}{29}\selectfont\color{covernavy}
     Otimização do Plano de Voo do Receptor em SAR Biestático};

  % ── subtítulo ────────────────────────────────────────────────────────────
  \node[anchor=north west, text width=12.6cm, align=left]
    at ([xshift=6.6cm,yshift=-14.2cm]current page.north west)
    {\coverfont\itshape\Large\color{coversub}
     Tomografia de Subsuperfície: Fundamento Físico,\\ Linha de Raciocínio e Utilidade Prática};

  % ── linha separadora ─────────────────────────────────────────────────────
  \draw[coverteal, line width=1.1pt]
    ([xshift=6.6cm,yshift=-16.9cm]current page.north west) --
    ++(5.4cm,0);

  % ── autor / projeto / data ──────────────────────────────────────────────
  \node[anchor=north west, text width=12.6cm, align=left]
    at ([xshift=6.6cm,yshift=-17.9cm]current page.north west)
    {\coverfont\bfseries\large\color{covernavy} David Atencia\\[0.35em]
     \coverfont\color{coversub}\normalsize Projeto de Pesquisa SAR-Tomo\\[0.15em]
     \coverfont\color{coversub}\normalsize \today};

\end{tikzpicture}
\end{titlepage}
\restoregeometry
\clearpage
\pagenumbering{arabic}

\begin{abstract}
\noindent
Este documento apresenta a proposta de otimização da trajetória do receptor em um sistema de
radar de abertura sintética (SAR) biestático voltado à tomografia de alvos subsuperficiais.
O transmissor percorre uma trajetória fixa em hélice cônica, enquanto a posição horizontal do
receptor é escolhida, a cada instante, de modo a maximizar a potência recebida — explorando o
\emph{ângulo de Brewster} da interface entre o ar e o solo — e, simultaneamente, minimizar a
resolução espacial tridimensional esperada do sistema tomográfico. Discute-se o caso de um meio
homogêneo de duas camadas (ar-solo) e sua generalização a um solo estratificado em várias camadas
planas, mais próximo de condições reais de campo. Apresentam-se o fundamento físico-matemático da
proposta, a linha de raciocínio que leva à estratégia de otimização adotada, os parâmetros de
projeto relevantes e os benefícios práticos esperados, incluindo a possibilidade de visualizar e
validar o plano de voo de forma georreferenciada antes da execução em campo.
\end{abstract}

\begin{nota}{Reprodutibilidade}
Os estudos, simulações e o código-fonte que fundamentam esta proposta encontram-se disponíveis
no repositório público do projeto:\\
\url{https://github.com/DavidAtenciaMondragon/sar_tomo_research}.
\end{nota}

\clearpage
\tableofcontents
\clearpage

%%═══════════════════════════════════════════════════════════════════════════════
\section{Introdução e Motivação}
%%═══════════════════════════════════════════════════════════════════════════════

Em um sistema SAR biestático para detecção de alvos enterrados, o transmissor percorre uma
trajetória predeterminada — uma \textbf{hélice cônica} — projetada para maximizar a diversidade
angular e, com isso, a resolução tridimensional do volume de interesse. A posição do receptor,
por outro lado, é um \textbf{grau de liberdade livre}: em cada instante de aquisição, pode ser
escolhida entre um número praticamente ilimitado de posições horizontais, mantendo a mesma
altitude do transmissor naquele instante.

A proposta aqui descrita responde à seguinte pergunta de projeto:

\begin{center}
\fbox{\parbox{0.85\textwidth}{\centering\itshape
Onde deve voar o receptor em cada instante para que o par biestático transmissor--receptor seja,
simultaneamente, ótimo em potência recebida \textbf{e} em resolução 3D do alvo subsuperficial?
}}
\end{center}

A proposta resolve esse problema de forma sistemática, posição a posição do transmissor ao longo
da hélice, produzindo uma trajetória ótima para o receptor que pode ser exportada de forma
georreferenciada e integrada a sistemas de informação geográfica (SIG) para planejamento e
validação do voo em campo.

\subsection{Por que o receptor precisa de um plano de voo próprio}

Em um radar monostático, transmissor e receptor compartilham a mesma plataforma e, portanto, a
mesma trajetória. Em um radar biestático essa restrição desaparece: o receptor pode ocupar uma
posição independente, escolhida propositalmente para melhorar o desempenho do sistema. Duas
oportunidades decorrem disso:

\begin{itemize}
  \item \textbf{Ganho de potência via ângulo de Brewster.} A interface ar-solo, para incidência
  em polarização TM, possui um ângulo específico ($\thetaB$) no qual a reflexão é nula e toda a
  potência incidente é transmitida ao subsolo. Posicionar o receptor de modo que o trecho
  alvo$\to$receptor incida próximo a esse ângulo aumenta a potência que efetivamente retorna.
  \item \textbf{Ganho de resolução via diversidade angular.} A resolução 3D de um sistema SAR
  biestático depende do ângulo de separação biestática entre transmissor e receptor, observado a
  partir do alvo. Escolher esse ângulo adequadamente reduz a área de resolução horizontal e
  vertical do voxel tomográfico.
\end{itemize}

Como se mostra na Seção~\ref{sec:custo}, esses dois objetivos nem sempre apontam para a mesma
posição do receptor, o que justifica uma formulação multiobjetivo e uma busca numérica — e não uma
fórmula fechada única — para cada posição do transmissor.

%%═══════════════════════════════════════════════════════════════════════════════
\section{Fundamentação Teórica}
%%═══════════════════════════════════════════════════════════════════════════════

\subsection{Geometria do sistema}

A Figura~\ref{fig:geometria} ilustra a configuração biestática. O transmissor percorre uma hélice
cônica no ar, entre um raio mínimo e máximo e uma faixa de altitudes. O volume de alvos
subsuperficiais é representado por um conjunto de pontos distribuídos no espaço 3D de interesse.
Para cada posição instantânea do transmissor, busca-se a posição horizontal do receptor — na mesma
altitude do transmissor naquele instante — que minimiza um critério combinado de desempenho.

\begin{figure}[H]
\centering
\begin{tikzpicture}[scale=0.95, every node/.style={font=\small}]
  % interface ar-solo
  \fill[pattern=north east lines, pattern color=brown!60] (-6,-2) rectangle (6,-3.6);
  \draw[thick] (-6,-2) -- (6,-2) node[right, black] {$z=0$ (interface ar--solo)};
  \node at (0,-2.9) {\textcolor{brown!70!black}{solo ($n_2$, ou pilha de camadas $n_1,\dots,n_M$)}};

  % Tx e helice
  \draw[-{Latex[length=2mm]}, blue!60!black, thick] (-3,3.2) .. controls (-1,2.6) and (1,2.2) .. (3.2,1.6);
  \filldraw[blue!60!black] (-3,3.2) circle (2pt) node[above] {Transmissor (início hélice)};
  \filldraw[blue!60!black] (3.2,1.6) circle (2.4pt) node[right] {Tx$_i$};

  % alvo
  \filldraw[red!70!black] (0,-2.6) circle (2pt) node[below] {alvo $P$};

  % Rx otimizado
  \filldraw[orange!80!black] (-2.6,1.4) circle (2.4pt) node[left] {Rx$_i$ (otimizado)};

  % raios refratados Tx->alvo
  \draw[blue!50!black, -{Latex[length=1.8mm]}] (3.2,1.6) -- (0.9,-2.0);
  \draw[blue!50!black, dashed, -{Latex[length=1.8mm]}] (0.9,-2.0) -- (0,-2.6);

  % raios refratados alvo->Rx
  \draw[orange!70!black, dashed, -{Latex[length=1.8mm]}] (0,-2.6) -- (-0.7,-1.95);
  \draw[orange!70!black, -{Latex[length=1.8mm]}] (-0.7,-1.95) -- (-2.6,1.4);

  \node[blue!50!black] at (2.4,0.15) {$R_T=R_1+R_2$};
  \node[orange!70!black] at (-1.85,-0.35) {$R_R=R_3+R_4$};

  \draw[<->, thick] (-2.6,1.55) .. controls (0,2.3) .. (3.2,1.75);
  \node at (0.2,2.55) {$\Delta\phi$ (ângulo biestático)};
\end{tikzpicture}
\caption{Geometria biestática: o transmissor percorre a hélice cônica; o receptor, em cada
instante, ocupa a posição $\mathrm{Rx}_i$ que otimiza conjuntamente a potência recebida (via
ângulo de Brewster) e a resolução 3D (via ângulo biestático $\Delta\phi$). Linhas contínuas: trecho
no ar; linhas tracejadas: trecho refratado no solo.}
\label{fig:geometria}
\end{figure}

\subsection{Interface ar-solo, polarização TM e ângulo de Brewster}
\label{sec:brewster}

O caso base considera dois meios homogêneos: ar ($z>0$, índice $\nuno=1$) e solo ($z<0$, índice
$\ndos=\sqrt{\varepsilon_r}$). A onda incide em polarização \textbf{TM} (campo elétrico no plano
de incidência), para a qual os coeficientes de Fresnel são
\[
r_{TM} = \frac{\ndos\cos\theta_1 - \nuno\cos\theta_2}{\ndos\cos\theta_1 + \nuno\cos\theta_2}, \qquad
t_{TM} = \frac{2\nuno\cos\theta_1}{\ndos\cos\theta_1 + \nuno\cos\theta_2},
\]
com $\theta_1,\theta_2$ relacionados pela lei de Snell ($\nuno\sin\theta_1=\ndos\sin\theta_2$).
Existe um ângulo de incidência particular — o \textbf{ângulo de Brewster} — para o qual $r_{TM}=0$
e, portanto, toda a potência incidente é transmitida ao segundo meio:
\begin{resultado}{Ângulo de Brewster}
\[
\boxed{\thetaB = \arctan\!\left(\frac{\ndos}{\nuno}\right)}
\]
Para um solo seco típico ($\varepsilon_r\approx 4 \Rightarrow \ndos=2$): $\thetaB\approx
63{,}43^\circ$, com ganho de transmitância de $\sim$12--14\% em relação a incidências
intermediárias típicas.
\end{resultado}

Essa é a base física que justifica orientar a posição do receptor: se os trechos
transmissor$\to$alvo e alvo$\to$receptor incidem na interface em ângulos próximos a $\thetaB$, a
potência que efetivamente atinge o receptor é maximizada.

\subsection{Princípio de Fermat e cálculo do caminho refratado}

O caminho físico entre uma fonte (transmissor ou receptor) e um alvo enterrado não é uma linha
reta: refrata-se na interface, segundo o princípio de Fermat (o caminho percorrido minimiza o
tempo de trânsito óptico $\mathcal{T}=\int n(\bm r)\,ds$). Para uma interface plana horizontal, o
ponto de refração ótimo pertence à reta que une as projeções horizontais da fonte e do alvo, o que
reduz um problema de otimização bidimensional a um problema unidimensional com solução única,
resolvido numericamente em poucas iterações. Esse resultado é o que torna viável avaliar um grande
número de posições candidatas do receptor sem custo computacional proibitivo.

\subsection{Equação de radar biestática com refração}

Incorporando a transmitância $T_1,T_2$ nas duas travessias da interface e os caminhos ópticos
totais $\RT=R_1+R_2$ (transmissor$\to$alvo) e $\RR=R_3+R_4$ (alvo$\to$receptor):
\begin{resultado}{Equação de radar biestática com refração}
\[
\boxed{\Prcv = \frac{\Pt\,G_T\,G_R\,\sigma\,T_1\,T_2\,\lambda^2}{(4\pi)^3\,\RT^2\,\RR^2}}
\]
onde $\Pt$ é a potência transmitida, $G_T,G_R$ são os ganhos de antena na direção do alvo, $\sigma$
é a seção reta radar do alvo e $\lambda$ é o comprimento de onda.
\end{resultado}
A dependência de $T_1$ e $T_2$ no ângulo de incidência é precisamente o elo entre a Seção
\ref{sec:brewster} e a potência recebida: maximizar $T_1T_2$ equivale a aproximar os ângulos de
incidência do ângulo de Brewster (ou de seu equivalente multicamada, Seção~\ref{sec:multicamada}).

\subsection{Resolução 3D biestática}

A resolução espacial do sistema tomográfico depende da geometria transmissor--alvo--receptor
através do ângulo biestático $\Delta\phi$ e do ângulo de observação. Em termos gerais, a resolução
horizontal $\dxy$ melhora (diminui) quando $\Delta\phi$ se afasta de $180^\circ$, e a resolução
vertical $\dz$ depende da largura de banda do sinal e da abertura sintética efetiva projetada ao
longo da hélice. Essas duas grandezas, calculadas para cada posição candidata do receptor, formam
o segundo critério de desempenho da proposta, ao lado da potência recebida.

\subsection{Extensão a solo estratificado (multicamada)}
\label{sec:multicamada}

Solos reais raramente são homogêneos: é comum encontrar camadas com propriedades dielétricas
distintas (por exemplo, solo seco superficial sobre solo mais úmido em profundidade). A proposta
generaliza o modelo de propagação a uma pilha de várias camadas planas e horizontais, cada uma
caracterizada por um índice de refração e uma espessura próprios.

\subsubsection{Invariante de Snell generalizado}
Para um raio que atravessa a pilha de camadas, o parâmetro do raio
\[
p = n_0\sin\theta_0 = n_1\sin\theta_1 = \cdots = \text{constante}
\]
é conservado em todas as interfaces. A distância horizontal total percorrida pelo raio é ajustada
numericamente até coincidir com a distância horizontal entre a fonte e o alvo, generalizando o
procedimento unidimensional de Fermat da Seção anterior para múltiplas interfaces.

\subsubsection{Transmitância multicamada}
Os coeficientes de reflexão e transmissão TM da pilha completa são obtidos por um formalismo de
matrizes de transferência que combina as contribuições de cada camada e se reduz exatamente às
equações de Fresnel de duas camadas (Seção~\ref{sec:brewster}) quando existe uma única interface.

\subsubsection{Ângulo ótimo de incidência: sem solução fechada}
Para uma única interface, o ângulo que maximiza a transmitância é o Brewster $\thetaB=\arctan(n_2/n_1)$.
Para uma pilha de várias camadas, efeitos de interferência entre interfaces sucessivas deslocam o
máximo de transmitância e não existe, em geral, solução analítica fechada. Por isso, o ângulo
ótimo é obtido por varredura numérica:
\begin{resultado}{Ângulo ótimo multicamada (numérico)}
\[
\theta_{\mathrm{opt}} = \operatorname*{arg\,max}_{\theta\in[0,\,\pi/2)} T_{TM}(\theta)
\]
calculado a partir das camadas entre a superfície e a camada que contém o alvo, e usado no lugar
de $\thetaB$ para orientar a busca da posição ótima do receptor.
\end{resultado}

\subsubsection{Reciprocidade}
Para meios sem perdas (índices reais), a transmitância de ida (transmissor$\to$alvo) é igual à de
volta (alvo$\to$receptor) para o mesmo parâmetro de raio, o que decorre da simetria do formalismo
multicamada e da conservação de energia — propriedade que reduz o custo computacional da busca ao
evitar recalcular quantidades já conhecidas do trecho fixo transmissor$\to$alvo.

%%═══════════════════════════════════════════════════════════════════════════════
\section{Linha de Raciocínio da Otimização}
%%═══════════════════════════════════════════════════════════════════════════════

\subsection{Por que combinar potência e resolução em um único critério}
\label{sec:custo}

Maximizar a potência recebida $\Prcv$ (tipicamente $10^{-10}$--$10^{-8}\,\text{W}$) e minimizar a
área de resolução $\dxy\cdot\dz$ (tipicamente $10^{-2}$--$10^{-4}\,\text{m}^2$) não são, em geral,
objetivos alcançados pela mesma posição do receptor. Uma combinação linear direta é problemática:
as unidades físicas são incompatíveis (W vs.\ m$^2$) e as escalas numéricas diferem em várias
ordens de grandeza, tornando qualquer peso linear arbitrário e sem interpretação física.

A solução adotada explora a natureza \emph{multiplicativa} da equação de radar e das fórmulas de
resolução, convertendo produtos em somas por meio do logaritmo:
\begin{resultado}{Critério multiobjetivo log-normalizado}
\[
\boxed{\Jcost = -(1-\alphaR)\log_{10}(\Prcv) + \alphaR\log_{10}(\dxy\cdot\dz)}
\]
\begin{itemize}[nosep]
  \item $\alphaR=0$: a busca considera \textbf{apenas a potência} recebida;
  \item $\alphaR=1$: a busca considera \textbf{apenas a resolução} 3D esperada;
  \item $\alphaR\in(0,1)$: equilíbrio multiobjetivo, com o peso relativo entre os dois critérios
  ajustável conforme a prioridade da campanha (valor de referência típico: $\alphaR=0{,}3$, ou
  seja, 70\% de peso em potência e 30\% em resolução).
\end{itemize}
A posição do receptor que minimiza $\Jcost$ é aquela que melhor equilibra potência recebida
elevada e resolução espacial fina, segundo o peso escolhido.
\end{resultado}

\subsection{Estratégia de busca com múltiplos pontos de partida}

O critério $\Jcost$, em função da posição horizontal do receptor, não é convexo: o ângulo
biestático e os coeficientes de reflexão introduzem múltiplos mínimos locais. Para reduzir o risco
de convergir a um mínimo local distante do ótimo global, a busca numérica parte de \textbf{vários
pontos de partida heurísticos}, escolhidos com base na geometria do problema, cobrindo:

\begin{enumerate}[nosep]
  \item a direção oposta ao transmissor em relação ao centro do volume de alvos (ângulo biestático
  próximo de $180^\circ$);
  \item duas direções perpendiculares à linha transmissor--centro do volume;
  \item posições guiadas pelo ângulo de Brewster (ou pelo ângulo ótimo multicamada), a uma
  distância horizontal característica do centro do volume, com pequena perturbação para evitar
  simetrias degeneradas.
\end{enumerate}
Cada ponto de partida é refinado por um método numérico de otimização com restrições de região
(a busca é limitada a um raio horizontal máximo em torno do transmissor), e a solução de menor
custo entre todos os pontos de partida é adotada como a posição ótima do receptor naquele
instante.

\subsection{Reaproveitamento do trecho fixo transmissor\texorpdfstring{$\to$}{->}alvo}

Como o transmissor e os alvos não mudam durante a busca da posição do receptor, todo o trecho
transmissor$\to$alvo (ponto de refração, ângulo de incidência, transmitância e distância óptica) é
calculado uma única vez por posição do transmissor, antes de avaliar as posições candidatas do
receptor. Apenas o trecho alvo$\to$receptor precisa ser recalculado a cada candidato avaliado, o
que reduz significativamente o esforço computacional total da busca.

\subsection{Resumo da linha de raciocínio}

\begin{resultado}{Resumo do procedimento (ambos os cenários — solo homogêneo e estratificado)}
Para cada posição do transmissor ao longo da hélice:
\begin{enumerate}[nosep]
  \item Calcular os ganhos de antena do transmissor em direção a cada alvo do volume de interesse.
  \item Calcular, uma única vez, o trecho fixo transmissor$\to$alvo (refração, transmitância,
  distância óptica).
  \item Definir a geometria de busca (ângulo de Brewster ou ângulo ótimo multicamada, distância
  característica associada).
  \item A partir de múltiplos pontos de partida, minimizar o critério combinado de potência e
  resolução, respeitando os limites da região de busca.
  \item Selecionar o ponto de menor custo como a posição ótima do receptor.
  \item Registrar a potência recebida, a resolução esperada e o valor do critério no ponto ótimo.
\end{enumerate}
Ao final, a trajetória completa do receptor é obtida por interpolação suave sobre as posições
ótimas calculadas, resultando em uma trajetória contínua e fisicamente executável, coerente com a
trajetória do transmissor.
\end{resultado}

%%═══════════════════════════════════════════════════════════════════════════════
\section{Parâmetros de Configuração}
%%═══════════════════════════════════════════════════════════════════════════════

Os parâmetros a seguir descrevem, do ponto de vista físico e de projeto, as grandezas que definem
cada cenário de voo. Eles podem ser ajustados conforme as condições reais do terreno e os
objetivos da campanha, sem alterar o fundamento teórico da proposta.

\subsection{Trajetória do transmissor (hélice cônica)}

\begin{table}[H]
\centering
\caption{Parâmetros geométricos e de operação do transmissor}
\label{tab:radar_params}
\small
\begin{tabular}{@{}p{4.6cm}p{6.6cm}p{2.8cm}@{}}
\toprule
Parâmetro & Papel no sistema & Valor de referência \\
\midrule
Raios mínimo e máximo da hélice & Definem a abertura angular vista pelo volume de alvos & $147{,}5$ / $172{,}5\,\text{m}$ \\
Altitudes mínima e máxima & Definem a inclinação da hélice cônica & $80$ / $120\,\text{m}$ \\
Número de voltas & Densidade de amostragem angular da trajetória & $2$ \\
Velocidade de voo & Determina o espaçamento entre pulsos ao longo da hélice & $120\,\text{m/s}$ \\
Frequência de repetição de pulso & Taxa de amostragem temporal do sinal & $400\,\text{Hz}$ \\
Frequência portadora & Define a penetração no solo e o comprimento de onda & $400\,\text{MHz}$ (banda P) \\
Largura de banda do sinal & Determina a resolução vertical (em profundidade) & $50\,\text{MHz}$ \\
Potência transmitida & Determina o alcance do enlace de rádio & $100\,\text{kW}$ \\
Ganho e abertura da antena & Concentram a energia na direção do volume de interesse & $6\,\text{dBi}$; aberturas de $14^\circ$ (azimute) e $50^\circ$ (elevação) \\
\bottomrule
\end{tabular}
\normalsize
\end{table}

A frequência portadora em banda P ($400\,\text{MHz}$, $\lambda\approx0{,}75\,\text{m}$) é escolhida
por seu baixo coeficiente de atenuação no solo ($\lesssim 0{,}5\,\text{dB/m}$ em solo seco), o que
permite detectar alvos a $5$--$10\,\text{m}$ de profundidade — inviável em bandas mais altas (por
exemplo, banda X), onde a atenuação supera $25\,\text{dB/m}$.

\subsection{Modelo do solo}

\begin{table}[H]
\centering
\caption{Parâmetros do meio subsuperficial nos dois cenários considerados}
\label{tab:soil_params}
\small
\begin{tabular}{@{}p{4.6cm}p{6.6cm}p{2.8cm}@{}}
\toprule
Cenário & Descrição & Valor de referência \\
\midrule
Solo homogêneo & Um único índice de refração para todo o subsolo & $n_2=2$ ($\varepsilon_r\approx4$) \\
Solo estratificado & Sequência de camadas planas, cada uma com índice de refração e espessura próprios, da superfície até a profundidade dos alvos & duas camadas de exemplo: $n{=}2$ ($1\,\text{m}$) sobre $n{=}3$ ($1\,\text{m}$) \\
\bottomrule
\end{tabular}
\normalsize
\end{table}

\subsection{Volume de alvos e critério de busca}

\begin{table}[H]
\centering
\caption{Parâmetros do volume de interesse e da busca da posição do receptor}
\label{tab:opt_params}
\small
\begin{tabular}{@{}p{4.6cm}p{6.6cm}p{2.8cm}@{}}
\toprule
Parâmetro & Papel no sistema & Valor de referência \\
\midrule
Extensão e profundidade do volume de alvos & Define a região subsuperficial de interesse tomográfico & dezenas de metros em planta; poucos metros de profundidade \\
Raio de busca do receptor & Limita a região horizontal onde o receptor pode ser posicionado & $200\,\text{m}$ em torno do transmissor \\
Peso potência--resolução ($\alphaR$) & Equilibra a prioridade entre potência recebida e resolução espacial & $0$ a $1$; valor de referência $0{,}3$ \\
Número de pontos de partida da busca & Reduz o risco de mínimos locais na otimização & tipicamente 5 \\
\bottomrule
\end{tabular}
\normalsize
\end{table}

\begin{nota}{Estudos de sensibilidade}
O peso $\alphaR$ e o tamanho do volume de alvos podem ser variados sistematicamente para
quantificar seu efeito na potência recebida, na resolução esperada e no critério combinado —
permitindo escolher, de forma fundamentada, a configuração mais adequada a cada campanha.
\end{nota}

\subsection{Resultados e uso em sistemas de informação geográfica}

Cada execução da proposta produz, para a trajetória completa do transmissor, a trajetória ótima
correspondente do receptor, acompanhada da potência recebida esperada, da resolução espacial
esperada e do valor do critério combinado ao longo do voo.

Como a geometria do problema é definida em relação a um ponto de referência geográfico conhecido,
os resultados podem ser exportados em formatos tabulares e georreferenciados, prontos para
importação em um sistema de informação geográfica (SIG), como o \textbf{QGIS}. Isso permite:
\begin{itemize}[nosep]
  \item visualizar, sobre imagens de satélite ou cartografia local, a trajetória planejada do
  transmissor e a trajetória otimizada do receptor;
  \item conferir a viabilidade operacional do voo do receptor (obstáculos, limites de propriedade,
  espaço aéreo) antes da campanha de campo;
  \item disponibilizar as coordenadas georreferenciadas para o piloto ou sistema de navegação do
  receptor.
\end{itemize}

%%═══════════════════════════════════════════════════════════════════════════════
\section{Benefícios da Proposta}
%%═══════════════════════════════════════════════════════════════════════════════

\begin{enumerate}
  \item \textbf{Aumento da potência recebida.} Ao orientar a posição do receptor para o ângulo de
  Brewster (ou seu equivalente numérico multicamada), a transmitância na interface se aproxima do
  máximo teórico, com ganho estimado de $\sim$12--14\% em relação a uma incidência típica
  arbitrária — relevante em um enlace de subsuperfície já limitado em potência.

  \item \textbf{Melhoria da resolução 3D tomográfica.} A escolha simultânea do ângulo biestático
  reduz a área de resolução do voxel de imageamento, sem exigir mudanças na trajetória (fixa) do
  transmissor.

  \item \textbf{Controle explícito do compromisso potência-resolução.} Um único peso ajustável
  permite priorizar sensibilidade de detecção (potência) ou nitidez de imageamento (resolução)
  conforme o objetivo da campanha, sem redesenhar a estratégia de busca.

  \item \textbf{Generalização a solos realistas.} A extensão a solo estratificado elimina a
  suposição de solo homogêneo, permitindo modelar estratificações comuns em campo (por exemplo,
  camada seca superficial sobre camada mais úmida) e obter um plano de voo consistente com a
  física real do terreno, ainda que sem perdas dielétricas modeladas (Seção~\ref{sec:limitacoes}).

  \item \textbf{Viabilidade computacional.} A redução do problema de refração a uma busca
  unidimensional, o reaproveitamento do trecho fixo transmissor$\to$alvo e a amostragem adequada
  da trajetória do transmissor tornam viável otimizar centenas de posições do receptor, com
  múltiplos pontos de partida cada, em tempo compatível com o planejamento operacional de uma
  campanha.

  \item \textbf{Integração direta com o fluxo de trabalho de campo.} A exportação georreferenciada
  dos resultados evita retrabalho manual de conversão de coordenadas e permite validar visualmente
  o plano de voo antes da execução, reduzindo o risco de replanejamento em campo.
\end{enumerate}

%%═══════════════════════════════════════════════════════════════════════════════
\section{Limitações e Trabalhos Futuros}
\label{sec:limitacoes}
%%═══════════════════════════════════════════════════════════════════════════════

\begin{itemize}[nosep]
  \item O modelo de propagação é \textbf{sem perdas} (índices de refração reais); solos reais têm
  permissividade complexa, que introduz atenuação adicional não modelada nesta proposta.
  \item As interfaces são consideradas \textbf{planas e horizontais}; rugosidade de superfície e
  heterogeneidade lateral do solo não são consideradas.
  \item O volume de alvos é representado por uma \textbf{amostragem discreta}, representativa do
  volume de interesse, não uma discretização exaustiva do meio contínuo real.
  \item A não convexidade do critério combinado é mitigada, mas não eliminada, pela busca com
  múltiplos pontos de partida; um número maior de pontos pode ser necessário em geometrias muito
  assimétricas.
  \item Trabalhos futuros: incorporar atenuação dielétrica complexa por camada e validar o plano de
  voo otimizado contra medições de campo.
\end{itemize}

%%═══════════════════════════════════════════════════════════════════════════════
\section*{Referências}
\addcontentsline{toc}{section}{Referências}
%%═══════════════════════════════════════════════════════════════════════════════

\begin{enumerate}[nosep]
  \item M.\ Born e E.\ Wolf, \emph{Principles of Optics}, 7ª ed., Cambridge University Press, 1999.
  \item F.\ T.\ Ulaby, R.\ K.\ Moore, A.\ K.\ Fung, \emph{Microwave Remote Sensing}, Artech House, 1981.
  \item A.\ Ishimaru, \emph{Electromagnetic Wave Propagation, Radiation, and Scattering}, Prentice Hall, 1991.
  \item J.\ A.\ Góes, \emph{Techniques for High-Resolution 3D Images with Synthetic Aperture Radar}, Tese de Doutorado, UNICAMP, 2022.
\end{enumerate}

\end{document}
